% AI-assisted manuscript source; responsible author: Carptopus.
\documentclass[11pt]{article}
\usepackage[a4paper,margin=29mm]{geometry}
\usepackage[T1]{fontenc}
\usepackage[utf8]{inputenc}
\usepackage{lmodern}
\usepackage{microtype}
\usepackage{amsmath,amssymb,amsthm,mathtools}
\usepackage{enumitem}
\usepackage{xcolor}
\usepackage[colorlinks=true,linkcolor=blue!55!black,citecolor=blue!55!black,urlcolor=blue!65!black]{hyperref}
\usepackage{fancyhdr}

\newtheorem{theorem}{Theorem}[section]
\newtheorem{lemma}[theorem]{Lemma}
\newtheorem{proposition}[theorem]{Proposition}
\theoremstyle{remark}
\newtheorem{remark}[theorem]{Remark}
\newtheorem*{maintheorem}{Theorem A}
\numberwithin{equation}{section}

\pagestyle{fancy}
\fancyhf{}
\fancyhead[L]{Carptopus}
\fancyhead[R]{Low-component total cut complexes}
\fancyfoot[C]{\thepage}
\setlength{\headheight}{14pt}
\setlength{\parindent}{0pt}
\setlength{\parskip}{0.55em}
\setlength{\emergencystretch}{4em}
\setlist{nosep,leftmargin=*}

\title{The low-component cases for total cut complexes\\of disconnected graphs}
\author{Carptopus\\\href{mailto:carptopus@163.com}{carptopus@163.com}}
\date{22 August 2026}

\hypersetup{
  pdftitle={The low-component cases for total cut complexes of disconnected graphs},
  pdfauthor={Carptopus},
  pdfsubject={A wedge-of-spheres formula for the low-component cases of total cut complexes},
  pdfkeywords={total cut complex, disconnected graph, Alexander duality, simplicial complex, wedge of spheres, polyhedral join}
}

\begin{document}
\maketitle

\begin{center}
\small\itshape
Public Beta v0.2. The proof, novelty boundary, and complete manuscript have passed
writing-admission, manuscript-level, and release-level zero-trust audits. External
mathematical review and journal peer review are pending.
\end{center}
\vspace{0.8em}

\begin{abstract}
Let $G=G_1\sqcup\cdots\sqcup G_k$ be a graph with $k$ nonempty connected
components and $n$ vertices. For $d\geq 2$, write $\Delta_d^t(G)$ for the total
$d$-cut complex and $\mathrm{BI}_d(G)$ for its Alexander dual. For $d\geq3$,
Carnero Bravo proved that, under natural contractibility and simple-connectivity
assumptions on the components, if $k\geq d+3$, then
\[
\Delta_d^t(G)\simeq
\bigvee_{\binom{k-1}{d-1}} S^{n-d-1},
\]
and asked whether the same formula holds for $k=d,d+1,d+2$. We also give the
elementary independent argument needed for the stated $d=2$ high-component range.

We answer all three cases affirmatively. For $k=d+2$, a complete $1$-skeleton
and a fixed-component triangular fan give simple connectivity. For $k=d+1$, we
separate the excess $n-k=0,1$ boundary cases and fill the fundamental cycles of a
spanning star by explicit triangular disks. For $k=d$, we cover the total cut
complex by joins indexed by weak compositions of $d$. The all-ones block is the
required sphere, while the union of the remaining blocks deformation retracts
onto its intersection with that sphere. The last argument requires only the
componentwise contractibility-or-void assumption and not the additional
simple-connectivity hypothesis. Together with the previously known range, this
completes the formula for every $k\geq d$ under the stated hypotheses.
\end{abstract}

\section{Introduction}

For a graph $G$ and an integer $d\geq1$, the total $d$-cut complex
$\Delta_d^t(G)$ consists of the vertex sets whose complement contains an
independent set of size $d$ \cite{ref1}. Equivalently,
\begin{equation}
\sigma\in\Delta_d^t(G)
\quad\Longleftrightarrow\quad
\alpha(G-\sigma)\geq d.
\label{eq:definition}
\end{equation}

Total cut complexes were introduced by Bayer, Denker, Jeli\'c Milutinovi\'c,
Rowlands, Sundaram, and Xue \cite{ref1}. Their topology connects graph structure,
Stanley--Reisner theory, Alexander duality, shellability, and discrete Morse
theory.

Carnero Bravo \cite{ref2} studies both $\Delta_d^t(G)$ and its Alexander dual
$\mathrm{BI}_d(G)$, the bounded-independence complex. For a disconnected graph
$G=G_1\sqcup\cdots\sqcup G_k$, \cite[Theorem 26]{ref2} calculates the
homotopy type of $\mathrm{BI}_d(G)$ under componentwise contractibility
assumptions. Alexander duality then determines the homology of $\Delta_d^t(G)$.
To upgrade this homology calculation to a wedge-of-spheres homotopy type,
\cite[Theorem 28]{ref2} proves simple connectivity by observing that every three
vertices form a face when $k\geq d+3$. Its proof invokes \cite[Theorem 26]{ref2},
which is stated for $d\geq3$; we therefore supply the short $d=2$ argument
independently in Section~\ref{sec:completion}.

This leaves exactly three component counts. In \cite[Question 29]{ref2}, it is
asked whether the same formula remains true for
\[
k=d,\qquad k=d+1,\qquad k=d+2.
\]

The difficulty is genuinely low-dimensional. The dual calculation already fixes
the target homology, but the complete-$2$-skeleton argument used for
$k\geq d+3$ is no longer available. The three missing layers exhibit three
different boundary phenomena:
\begin{enumerate}
\item for $k=d+2$, the $1$-skeleton is complete, but some triangles may be absent;
\item for $k=d+1$, even the $1$-skeleton can have missing edges and the target may
have dimension $0$ or $1$;
\item for $k=d$, the all-singleton case is already a sphere, and the extra faces
contributed by larger components must be shown not to change its homotopy type.
\end{enumerate}

Our main theorem resolves all three layers.

\begin{maintheorem}\label{thm:main}
Let $d\geq2$, and let
\[
G=G_1\sqcup\cdots\sqcup G_k
\]
have $k\geq d$ nonempty connected components and $n$ vertices. Assume that
\begin{enumerate}
\item for every $i$ and every $2\leq \ell\leq d$, the complex
$\Delta_\ell^t(G_i)$ is void or contractible; and
\item $\mathrm{BI}_2(G_i)$ is simply connected for every $i$.
\end{enumerate}
Then
\begin{equation}
\Delta_d^t(G)\simeq
\bigvee_{\binom{k-1}{d-1}} S^{n-d-1}.
\label{eq:main}
\end{equation}
For $d\geq3$, the range $k\geq d+3$ is \cite[Theorem 28]{ref2}. The main new
content is the range $k=d,d+1,d+2$; Section~\ref{sec:completion} also closes the
elementary $d=2$, $k\geq5$ dependency separately. Moreover, when $k=d$,
condition~2 is unnecessary.
\end{maintheorem}

The proof for $k=d$ is naturally expressed by a weak-composition cover. This is
close in spirit to the composition-poset and two-factor decompositions developed
in \cite[Theorems 26 and 34]{ref2} and to the general polyhedral-join framework
in \cite{ref3}. Those tools are prior work. The point here is the particular
multi-filter cover and the comparison between its nonbaseline union and its trace
on the baseline sphere.

\section{Definitions and conventions}

All graphs are finite and simple. All simplicial complexes are finite. We
distinguish the void complex, which has no faces, from $\{\varnothing\}$. A
contractible complex is in particular nonvoid.

For a graph $H$ and $r\geq1$, set
\begin{equation}
\Delta_r^t(H)
=\{\sigma\subseteq V(H):\alpha(H-\sigma)\geq r\}.
\label{eq:total-cut}
\end{equation}
The bounded-independence complex is
\begin{equation}
\mathrm{BI}_r(H)
=\{\sigma\subseteq V(H):\alpha(H[\sigma])<r\}.
\label{eq:bounded-independence}
\end{equation}
These two complexes are Alexander dual on the common vertex set \cite{ref2}. We
use the standard join $K*L$ for complexes on disjoint vertex sets.

For the weak-composition argument, it is convenient to extend the total-cut
filtration by
\begin{equation}
\Delta_0^t(H)=\Delta^{V(H)},
\qquad
\Delta_1^t(H)=\partial\Delta^{V(H)}.
\label{eq:extended-filtration}
\end{equation}
The second identity follows because $H-\sigma$ contains an independent vertex
exactly when $\sigma\neq V(H)$. If $H$ has one vertex, then
$\partial\Delta^{V(H)}=\{\varnothing\}$, viewed as $S^{-1}$ in join formulas.

We repeatedly use the following elementary observation.

\begin{lemma}\label{lem:components}
If deleting a vertex set $\sigma$ leaves at least $d$ nonempty connected
components of $G$, then $\sigma\in\Delta_d^t(G)$.
\end{lemma}

\begin{proof}
Choose one remaining vertex from each of $d$ components. Vertices in distinct
components are pairwise nonadjacent, so they form an independent $d$-set.
\end{proof}

\section{The inherited homology calculation}

The component hypotheses first determine the homology needed in the two
connectivity arguments.

\begin{proposition}\label{prop:homology}
Under the hypotheses of Theorem~A, $\Delta_d^t(G)$ has only one nonzero reduced
homology group, namely
\begin{equation}
\widetilde H_{n-d-1}(\Delta_d^t(G);\mathbb Z)
\cong \mathbb Z^{\binom{k-1}{d-1}}.
\label{eq:homology}
\end{equation}
\end{proposition}

\begin{proof}
Suppose first that $d\geq3$. The proof of \cite[Theorem 28]{ref2}, using its
connectivity result for bounded-independence complexes, shows that the component
hypotheses make every $\mathrm{BI}_\ell(G_i)$ contractible for
$2\leq\ell\leq d$. Then \cite[Theorem 26]{ref2} gives
\begin{equation}
\mathrm{BI}_d(G)\simeq
\bigvee_{\binom{k-1}{d-1}}S^{d-2}.
\label{eq:dual-wedge}
\end{equation}
Combinatorial Alexander duality gives \eqref{eq:homology}.

For $d=2$, $\mathrm{BI}_2(G)$ is the disjoint union, in geometric realization,
of the clique complexes $\mathrm{BI}_2(G_i)$. Each is contractible under the
assumptions: if the dual total-cut complex is void this is immediate, and if it
is contractible, Alexander duality gives acyclicity while the assumed simple
connectivity permits Whitehead's theorem. Thus $\mathrm{BI}_2(G)$ is homotopy
equivalent to $k$ points. Its reduced $H_0$ has rank $k-1$, which is
$\binom{k-1}{1}$, and Alexander duality again yields \eqref{eq:homology}.
\end{proof}

For $k=d+1,d+2$, we will combine Proposition~\ref{prop:homology} with direct
low-dimensional connectivity arguments. The $k=d$ proof is stronger and does
not use this proposition.

\section{The case \texorpdfstring{$k=d$}{k=d}}

We first treat the most structured boundary layer.

\begin{theorem}\label{thm:k-equals-d}
Let $d\geq2$, and let $G=G_1\sqcup\cdots\sqcup G_d$ have $d$ nonempty connected
components and $n$ vertices. If $\Delta_\ell^t(G_i)$ is void or contractible for
every $i$ and $2\leq\ell\leq d$, then
\begin{equation}
\Delta_d^t(G)\simeq S^{n-d-1}.
\label{eq:k-equals-d}
\end{equation}
\end{theorem}

\begin{proof}
Let
\[
W_d=\{(r_1,\ldots,r_d)\in\mathbb Z_{\geq0}^d:
r_1+\cdots+r_d=d\}
\]
be the weak compositions of $d$ into $d$ parts. For
$\mathbf r=(r_1,\ldots,r_d)\in W_d$, define
\begin{equation}
X_{\mathbf r}
=\Delta_{r_1}^t(G_1)*\cdots*\Delta_{r_d}^t(G_d).
\label{eq:blocks}
\end{equation}
We claim that
\begin{equation}
\Delta_d^t(G)=\bigcup_{\mathbf r\in W_d}X_{\mathbf r}.
\label{eq:block-cover}
\end{equation}

Indeed, let $\sigma\in\Delta_d^t(G)$ and choose an independent $d$-set
$S\subseteq V(G)-\sigma$. Put $r_i=|S\cap V(G_i)|$. Then
$\mathbf r\in W_d$, and the component of $\sigma$ in $V(G_i)$ lies in
$\Delta_{r_i}^t(G_i)$. Hence $\sigma\in X_{\mathbf r}$. Conversely, if
$\sigma\in X_{\mathbf r}$, choose an independent $r_i$-set in
each $G_i-(\sigma\cap V(G_i))$. Their union is an independent $d$-set in
$G-\sigma$. This proves \eqref{eq:block-cover}.

The unique weak composition with every coordinate equal to $1$ is
$\mathbf 1=(1,\ldots,1)$. Its block is
\begin{equation}
B=X_{\mathbf1}
=\partial\Delta^{V(G_1)}*\cdots*\partial\Delta^{V(G_d)}
\simeq S^{n-d-1}.
\label{eq:baseline}
\end{equation}
Let $R=W_d-\{\mathbf1\}$, omit any void blocks, and put
\begin{equation}
Y=\bigcup_{\mathbf r\in R}X_{\mathbf r},
\qquad
A=B\cap Y.
\label{eq:nonbaseline}
\end{equation}
Every $\mathbf r\in R$ has a zero coordinate and a coordinate at least $2$.
The zero coordinate supplies a full-simplex factor in \eqref{eq:blocks}, so
every nonvoid $X_{\mathbf r}$ is a cone.

Now take a finite nonempty set $T\subseteq R$ and define
\[
m_i=\max_{\mathbf r\in T}r_i.
\]
Because the total-cut filtration decreases with its index and the component
vertex sets are disjoint,
\begin{equation}
\bigcap_{\mathbf r\in T}X_{\mathbf r}
=\Delta_{m_1}^t(G_1)*\cdots*\Delta_{m_d}^t(G_d).
\label{eq:cover-intersection}
\end{equation}
At least one $m_i$ is at least $2$. If some $m_j=0$, the intersection is a cone.
If no coordinate is zero, every nonvoid factor of index at least $2$ is
contractible by hypothesis. Thus every nonempty intersection in this cover of
$Y$ is contractible.

Intersecting \eqref{eq:cover-intersection} with $B$ gives
\begin{equation}
B\cap\bigcap_{\mathbf r\in T}X_{\mathbf r}
=\Delta_{\max(1,m_1)}^t(G_1)*\cdots*
\Delta_{\max(1,m_d)}^t(G_d).
\label{eq:trace-intersection}
\end{equation}
Every nonempty complex in \eqref{eq:trace-intersection} is contractible for the
same reason. Moreover, \eqref{eq:cover-intersection} contains a nonempty simplex
if and only if \eqref{eq:trace-intersection} does. Only the zero coordinates
change, from a full simplex to a simplex boundary; a factor with index at least
$2$ is nonvoid and contractible, hence contains a nonempty face. Therefore the
covers of $Y$ and $A$ have exactly the same nonempty-intersection pattern.

Take as a basis-like cover of $Y$ all nonempty finite intersections of the
$X_{\mathbf r}$. For the inclusion $p:A\hookrightarrow Y$, the inverse image of
each cover element is its intersection with $B$. By
\eqref{eq:cover-intersection}--\eqref{eq:trace-intersection}, both spaces are
contractible. The local-to-global theorem \cite[Theorem 10]{ref2} implies that
$p$ is a weak equivalence, hence a homotopy equivalence between finite CW
complexes.

The inclusion $A\hookrightarrow Y$ is also a cofibration, so $A$ is a strong
deformation retract of $Y$. Extending this deformation by the identity on $B$
gives
\[
\Delta_d^t(G)=B\cup_A Y\simeq B\simeq S^{n-d-1}.
\]
If $R$ contains no nonvoid block, then $\Delta_d^t(G)=B$ from the outset.
\end{proof}

\begin{remark}
Theorem~\ref{thm:k-equals-d} does not use the simple connectivity of
$\mathrm{BI}_2(G_i)$. It is therefore slightly stronger than the $k=d$
specialization requested in \cite[Question 29]{ref2}.
\end{remark}

\section{The case \texorpdfstring{$k=d+1$}{k=d+1}}

Write
\begin{equation}
e=n-k=\sum_{i=1}^k(|V(G_i)|-1).
\label{eq:excess}
\end{equation}
The target sphere dimension is exactly $e$.

\begin{theorem}\label{thm:k-equals-d-plus-one}
Under the hypotheses of Theorem~A, if $k=d+1$, then
\begin{equation}
\Delta_d^t(G)\simeq\bigvee_d S^{n-d-1}.
\label{eq:k-equals-d-plus-one}
\end{equation}
\end{theorem}

\begin{proof}
If $e=0$, all components are singletons. Then
\[
\Delta_d^t(G)=\{\sigma:|\sigma|\leq1\},
\]
which is a set of $d+1$ points and hence a wedge of $d$ copies of $S^0$.

If $e=1$, there is one two-vertex connected component $\{a,b\}$ and $d$
singleton components. The complex is the graph obtained from $K_{2,d}$ by
adding the edge $ab$. It is connected and has cycle rank
\[
(2d+1)-(d+2)+1=d.
\]
Thus it is a wedge of $d$ circles.

Assume now that $e\geq2$. If a component with at least three vertices exists,
choose one such component as $C$; otherwise choose any nonsingleton component as
$C$, in which case every nonsingleton component is a $K_2$. Fix $a\in C$. For
every other vertex $x$, deleting $a$ and $x$ empties at most one component. By
Lemma~\ref{lem:components}, $ax$ is an edge of $\Delta_d^t(G)$. Hence the star
with center $a$ is a spanning tree of the $1$-skeleton. The fundamental group is
generated by loops
\begin{equation}
a-x-y-a,
\label{eq:fundamental-loop}
\end{equation}
where $xy$ is an edge not in the star.

First suppose $|C|\geq3$. Choose a neighbor $b$ of $a$ in the connected graph
$C$. If $axy$ is a face, the loop \eqref{eq:fundamental-loop} is filled.
Otherwise $b\notin\{x,y\}$. The triples $abx$ and $aby$ are faces, since
deleting either triple empties at most one component. Because $xy$ is an edge,
there is an independent $d$-set $S$ in $G-\{x,y\}$. The failure of $axy$ to be
a face forces $a\in S$. Since $ab$ is an edge of $G$, we have $b\notin S$, so
the same $S$ witnesses that $bxy$ is a face. The triangles $abx$, $aby$, and
$bxy$ form a disk with boundary \eqref{eq:fundamental-loop}.

It remains to consider the case in which every nonsingleton component is a
$K_2$. Since $e\geq2$, there are at least two such components. Write one as
$\{a,b\}$ and choose a vertex $c$ in another. If $axy$ is not a face while $xy$
is an edge, then necessarily $\{x,y\}=\{b,s\}$ for a singleton component $s$.
Indeed, if $b$ is absent, failure of $ab$'s star triangle requires $x$ and $y$
to exhaust two singleton components, contradicting that $xy$ is an edge; if
$b$ is present and the other point is not a singleton, at most one component is
exhausted. In the remaining case, $abc$, $acs$, and $bcs$ are faces and form a
disk with boundary $a-b-s-a$.

Every generator \eqref{eq:fundamental-loop} is therefore null-homotopic, so
$\Delta_d^t(G)$ is simply connected for $e\geq2$.
Proposition~\ref{prop:homology} says that its only nonzero reduced homology is
$\mathbb Z^d$ in dimension $e$. The standard wedge-of-spheres criterion used in
\cite[Theorem 5]{ref2} now gives \eqref{eq:k-equals-d-plus-one}.
\end{proof}

\section{The case \texorpdfstring{$k=d+2$}{k=d+2}}

\begin{theorem}\label{thm:k-equals-d-plus-two}
Under the hypotheses of Theorem~A, if $k=d+2$, then
\begin{equation}
\Delta_d^t(G)\simeq
\bigvee_{\binom{d+1}{2}}S^{n-d-1}.
\label{eq:k-equals-d-plus-two}
\end{equation}
\end{theorem}

\begin{proof}
Deleting any two vertices empties at most two components, so at least $d$
nonempty components remain. Lemma~\ref{lem:components} shows that the
$1$-skeleton of $\Delta_d^t(G)$ is complete.

Suppose some component $C$ has at least two vertices, and fix $v\in C$. For any
two other vertices $x,y$, the triple $vxy$ cannot exhaust three components:
exhausting a component with $v$ requires another vertex of $C$, leaving only one
vertex available to exhaust another component. Hence $vxy$ is a face. Every
loop in the complete $1$-skeleton is filled by a triangular fan with apex $v$.
Thus the complex is simply connected. Here $n\geq d+3$, so the target dimension
$n-d-1$ is at least $2$. Proposition~\ref{prop:homology} and the
wedge-of-spheres criterion give \eqref{eq:k-equals-d-plus-two}.

If every component is a singleton, then $G=(d+2)K_1$ and
\[
\Delta_d^t(G)=\{\sigma:|\sigma|\leq2\}
=\mathrm{sk}_1\Delta^{d+1}.
\]
This is the complete graph on $d+2$ vertices. Its cycle rank is
\[
\binom{d+2}{2}-(d+2)+1=\binom{d+1}{2},
\]
which again gives \eqref{eq:k-equals-d-plus-two}.
\end{proof}

\section{Completion of the disconnected formula}\label{sec:completion}

\begin{proof}[Proof of Theorem~A]
The cases $k=d$, $k=d+1$, and $k=d+2$ are
Theorems~\ref{thm:k-equals-d}, \ref{thm:k-equals-d-plus-one}, and
\ref{thm:k-equals-d-plus-two}. If $d\geq3$, the range $k\geq d+3$ is
\cite[Theorem 28]{ref2}.

It remains only to make the $d=2$, $k\geq5$ range independent of the parameter
restriction in \cite[Theorem 26]{ref2}. Proposition~\ref{prop:homology} gives
\[
\widetilde H_{n-3}(\Delta_2^t(G);\mathbb Z)
\cong\mathbb Z^{k-1}
\]
and no other reduced homology. Deleting any three vertices empties at most three
components, so at least $k-3\geq2$ nonempty components remain.
Lemma~\ref{lem:components} shows that the $2$-skeleton of $\Delta_2^t(G)$ is the
full $2$-skeleton on its vertex set. Hence the complex is simply connected.
Since $n-3\geq2$, the wedge-of-spheres criterion \cite[Theorem 5]{ref2} gives
\[
\Delta_2^t(G)\simeq\bigvee_{k-1}S^{n-3}.
\]
These ranges exhaust all $d\geq2$ and $k\geq d$.
\end{proof}

The binomial multiplicities agree on the three boundary layers:
\[
\binom{d-1}{d-1}=1,
\qquad
\binom{d}{d-1}=d,
\qquad
\binom{d+1}{d-1}=\binom{d+1}{2}.
\]

The proofs also explain why the single complete-$2$-skeleton argument of
\cite[Theorem 28]{ref2} does not simply extend downward. At $k=d+2$ only selected
triangles are needed; at $k=d+1$ one must first choose a spanning star and fill
its fundamental cycles; at $k=d$ the correct object is the entire
weak-composition cover rather than the $2$-skeleton.

\section{Relation to previous decomposition methods}

The proof uses several established tools whose scope should be separated from
the new conclusion.

First, for $d\geq3$, the homology input is inherited from the
bounded-independence calculation and Alexander duality in
\cite[Theorems 26 and 28]{ref2}. For $d=2$, Proposition~\ref{prop:homology}
supplies the elementary clique-component/Alexander-duality argument directly.
Second, the passage from local equivalences on a basis-like cover to a global
weak equivalence is \cite[Theorem 10]{ref2}. Third, \cite[Theorem 34]{ref2} gives
a two-graph decomposition of $\Delta_d^t(G_1\sqcup G_2)$ when the relevant
filtration inclusions are null-homotopic. This already identifies
independent-number splitting as a natural mechanism for disjoint unions.

Theorem~\ref{thm:k-equals-d} is not a direct iteration of that two-factor result.
After combining several components, the intermediate total-cut filtration need
not retain the componentwise contractibility and null-homotopy hypotheses. The
weak-composition cover instead handles all distributions simultaneously and
compares the full nonbaseline diagram to its trace on the all-ones block.

Polyhedral joins and their homotopy-colimit decompositions are also established
objects; see \cite{ref3}. We do not claim a new general theorem about polyhedral
joins. The retained point is the total-cut-specific cover and its consequence
for the three component counts asked about in \cite[Question 29]{ref2}.

\section{Boundary of the result}

Theorem~A is conditional on the same component assumptions as
\cite[Theorem 28]{ref2}. It does not classify $\Delta_d^t(G)$ for arbitrary
disconnected graphs. In particular, it does not answer \cite[Question 30]{ref2},
which asks whether the clique-complex connectivity condition can be removed.

The $k=d$ argument shows that one hypothesis can be removed in that single
layer, but it does not imply a corresponding weakening for $k>d$. Nor does the
proof determine simple-homotopy type, shellability, torsion outside the stated
hypotheses, or the maps between consecutive total-cut filtration levels.

No computer enumeration is used in the proof. Finite calculations performed
during discovery served only as falsification controls and are not evidence for
the universal statements.

\section{AI-assisted research and writing disclosure}

This work was developed with AI assistance for literature discovery, proof
stress-testing, symbolic organization, and drafting. Carptopus is the responsible
author and has reviewed the mathematical claims and the disclosed computational
boundaries. OpenAI Codex is a research and writing tool, not an author.

\begingroup
\small
\raggedright
\setlength{\parskip}{0.25em}
\begin{thebibliography}{9}

\bibitem{ref1}
M. Bayer, M. Denker, M. Jeli\'c Milutinovi\'c, R. Rowlands, S. Sundaram, and
L. Xue,
\emph{Total Cut Complexes of Graphs}, \emph{Discrete \& Computational Geometry}
(2024), published version.
\url{https://doi.org/10.1007/s00454-024-00630-4}

\bibitem{ref2}
A. Carnero Bravo,
\emph{Total cut complexes and their duals}, arXiv:2602.21427v3 (2026).
\url{https://arxiv.org/abs/2602.21427}

\bibitem{ref3}
A. Carnero Bravo,
\emph{Polyhedral joins and graph complexes}, arXiv:2307.14321v3 (2025).
\url{https://arxiv.org/abs/2307.14321}

\end{thebibliography}
\endgroup

\end{document}
